\documentclass[11pt]{article}

\usepackage[margin=1in]{geometry}
\usepackage{booktabs}
\usepackage{amsmath}
\usepackage{amssymb}
\usepackage{array}
\usepackage{hyperref}
\usepackage{xcolor}
\usepackage{enumitem}
\usepackage{longtable}

\hypersetup{
    colorlinks=true,
    linkcolor=blue,
    urlcolor=blue,
    citecolor=blue
}

\setlist[itemize]{leftmargin=1.2em}
\setlist[enumerate]{leftmargin=1.4em}

\title{Stage 1 Report: Recent-Window MLP for CSI500 3-Day Stock Selection}
\author{Aaron Gao}
\date{\today}

\begin{document}
\maketitle

\begin{abstract}
This report documents my Stage 1 stock-selection system for the CSI500
competition. The submitted model is a recent-window multi-layer perceptron
(MLP) trained to rank CSI500 constituents by 3-trading-day forward return and
converted into a constrained long-only portfolio. The main contribution is not
a new black-box architecture, but a sequence of controlled experiments showing
that for this short horizon, restricting the training sample to recent market
history improves the validation and held-out portfolio metrics more reliably
than adding generic model complexity. The final Stage 1 submission uses the
\texttt{mlp/day3\_stage1\_recent\_window} branch. In the canonical held-out
self-test, it achieved a mean 3-day excess return of \(+2.526\%\), compared
with \(+1.744\%\) for the previous tuned MLP and \(+1.394\%\) for the Stage 1
XGBoost model. Under a longer robust split it remained positive at
\(+0.804\%\), and under a 3-window walk-forward check it achieved a weighted
mean excess return of \(+0.783\%\).
\end{abstract}

\section{Introduction and Problem Formulation}

The Stage 1 task is to construct a long-only portfolio over CSI500 constituents
using public historical market data. For each submission date \(t\), the model
assigns portfolio weights \(w_{i,t}\) to selected stocks \(i\). The portfolio
must satisfy
\[
    w_{i,t} \ge 0,\qquad
    \sum_i w_{i,t} = 1,\qquad
    w_{i,t} \le 0.10,\qquad
    \#\{i:w_{i,t}>0\}\ge 30.
\]
The target horizon for Stage 1 is 3 trading days. For stock \(i\), I define
the raw forward return as
\[
    r_{i,t}^{(3)} = \frac{P_{i,t+3}}{P_{i,t}} - 1,
\]
where \(P_{i,t}\) is the adjusted close price. The realized portfolio return is
\[
    R_{p,t}^{(3)} = \sum_i w_{i,t} r_{i,t}^{(3)}.
\]
The benchmark is the CSI500 index return over the same horizon,
\[
    R_{b,t}^{(3)} = \frac{I_{t+3}}{I_t} - 1.
\]
The primary evaluation quantity used throughout my experiments is mean excess
return,
\[
    \overline{R}_{excess}
    =
    \frac{1}{T}\sum_{t=1}^{T}\left(R_{p,t}^{(3)}-R_{b,t}^{(3)}\right).
\]
I also monitor positive excess rate and rank information coefficient (rank IC).
Rank IC is computed as the average daily Spearman correlation between predicted
scores and realized 3-day returns. It measures whether the model ranks stocks
correctly even before applying portfolio constraints.

\section{Data and Factors}

\subsection{Data Sources}

The system uses public historical market data downloaded by the course-provided
script. For the final Stage 1 submission, I followed the competition guideline
and refreshed the dataset with
\[
    \texttt{python download\_data.py --update --end 20260503}.
\]
The update process refreshes stock bars, the CSI500 benchmark series, and the
constituent list. The final generated submission file used prediction date
\texttt{2026-04-30}. I did not use private data, pretrained models, or
non-public alternative data. All reported experiments use public price,
volume, amount, turnover when available, constituent information, and the
CSI500 index.

\subsection{Factor Design}

Across the project I deliberately kept the data source fixed and changed the
factor transformations and model class one step at a time. This makes the
experiments easier to interpret: if a later model improves, the improvement is
less likely to come from an undocumented data difference. All submitted and
tested systems use only public price, volume, turnover, constituent, and CSI500
index data. The main signal groups are:

\begin{itemize}
    \item \textbf{Recent returns and momentum:}
    \(1,2,3,5,10,20\)-day returns, lagged 1-day returns, and moving-average
    distance features. These are included because the Stage 1 holding horizon
    is short and recent price behavior is likely more relevant than long-run
    fundamentals.
    \item \textbf{Market-relative returns:}
    stock returns minus CSI500 index returns over multiple horizons, together
    with relative-strength features such as
    \texttt{relative\_strength\_3\_10}. These factors help separate
    stock-specific strength from broad market movement.
    \item \textbf{Intraday and range features:}
    intraday return, opening gap, daily range, range relative to the recent
    average, and close position within the daily range. These summarize short
    horizon supply-demand pressure.
    \item \textbf{Risk and volatility:}
    rolling volatility over 3, 5, 10, and 20 days, downside volatility,
    volatility ratio, 20-day beta, idiosyncratic volatility, and idiosyncratic
    momentum. These features are useful both for prediction and for avoiding
    portfolios dominated by unstable names.
    \item \textbf{Liquidity and trading activity:}
    standardized volume, standardized amount, turnover z-scores, and turnover
    moving averages. Liquidity information is important because extreme short
    horizon moves often coincide with abnormal trading activity.
    \item \textbf{Cross-sectional ranks:}
    rank-normalized versions of several momentum, volume, beta, volatility,
    and price-location features. Since the portfolio is selected cross
    sectionally each day, rank features are a natural way to make the model
    compare stocks within the same market environment.
\end{itemize}

Rows with missing features or missing 3-day target are excluded from supervised
training. This also prevents leakage from the last 3 trading days, whose
future target is not yet observed.

\subsection{Factors Used in Each Attempt}

Table \ref{tab:factor-by-attempt} summarizes the factor design for each major
attempt. The important point is that the final model did not win because it
used private data or a completely different information set. It won because the
same public-price factor set was paired with a training window that better
matched the short 3-day horizon.

{\small
\setlength{\tabcolsep}{4pt}
\begin{longtable}{p{0.18\textwidth}p{0.36\textwidth}p{0.36\textwidth}}
\caption{Signals used in each modeling attempt and the reason for including them.}
\label{tab:factor-by-attempt}\\
\toprule
Attempt & Signals used & Why these signals were reasonable \\
\midrule
\endfirsthead
\toprule
Attempt & Signals used & Why these signals were reasonable \\
\midrule
\endhead
3-day XGBoost baseline &
Basic technical factors: recent returns, moving-average distance, volatility,
volume and turnover activity, and simple cross-sectional ranks. &
This created a clean benchmark using only standard price-volume information.
It tested whether a nonlinear tree model could extract short-horizon signal
without additional feature engineering. \\
\midrule
Ridge enhanced &
Expanded 3-day factor set: short returns, market-relative returns, intraday
return, opening gap, daily range, downside volatility, beta, idiosyncratic
momentum, liquidity z-scores, and rank-normalized versions of these signals. &
Ridge is linear, so richer explicit transformations are necessary. These
factors directly target the competition horizon: recent price pressure,
stock-specific strength versus CSI500, liquidity shocks, and risk-adjusted
movement. \\
\midrule
Stage 1 XGBoost &
The same enhanced public-price factor panel as Ridge, with nonlinear tree
splits over momentum, volatility, liquidity, and market-relative variables. &
XGBoost was used to test whether interactions among the engineered factors
were useful, for example whether momentum only helped under certain volatility
or liquidity conditions. \\
\midrule
Base and tuned MLP &
The enhanced factor panel after standardization, including return, rank,
market-relative, liquidity, volatility, beta, and idiosyncratic components. &
The MLP requires scaled continuous inputs and can learn smoother nonlinear
interactions than tree splits. This attempt tested whether a compact neural
cross-sectional model improved over XGBoost. \\
\midrule
LSTM sequence model &
Recent multi-day sequences of the same market features, using a 20-day lookback
sequence for each stock. &
This tested whether explicitly modeling temporal order helped beyond using
rolling technical summaries. It was included because the target is short term,
but the available panel is not very long. \\
\midrule
Style-dynamic and regime branches &
The enhanced stock-level factors plus market-state summaries such as index
returns, realized volatility, drawdown, breadth, and trend indicators. &
These branches tested whether the best stock-selection or portfolio rule should
change between bullish, neutral, and stressed market states. \\
\midrule
Recent-window MLP &
The same enhanced MLP factor panel, but trained only on recent windows
\((63,126,189,252,\text{full})\) instead of always using full history. &
For a 3-day horizon, factor behavior can change quickly. Restricting the sample
is a way to keep the training distribution close to the submission regime while
still using enough cross-sectional observations. \\
\midrule
Excess-target / denoised variants &
The recent-window factor panel with labels transformed to excess return,
winsorized targets, z-scored targets, and recency sample weights. &
These variants tested whether aligning the supervised label more directly with
the final excess-return metric would help. They were plausible, but the held-out
results showed that the transformation removed useful raw-return signal. \\
\bottomrule
\end{longtable}
}

\section{Model and Portfolio Construction}

\subsection{Models Attempted}

I treated the project as a sequence of controlled model-building attempts
rather than a single final run. Each attempt changed one of three components:
the factor representation, the prediction model, or the portfolio mapping. The
attempts are summarized in Table \ref{tab:model-by-attempt}.

{\small
\setlength{\tabcolsep}{4pt}
\begin{longtable}{p{0.16\textwidth}p{0.32\textwidth}p{0.25\textwidth}p{0.17\textwidth}}
\caption{How each prediction model was trained and how its scores were converted to a portfolio.}
\label{tab:model-by-attempt}\\
\toprule
Attempt & Prediction model and training & Portfolio construction & Selection criterion \\
\midrule
\endfirsthead
\toprule
Attempt & Prediction model and training & Portfolio construction & Selection criterion \\
\midrule
\endhead
3-day XGBoost baseline &
Gradient-boosted trees trained to predict raw 3-day forward return. This was
the first nonlinear benchmark and used chronological validation. &
Top-ranked stocks were selected and converted to long-only weights under the
10\% cap. &
Validation mean excess return. \\
\midrule
Ridge enhanced &
Ridge regression trained on the expanded 3-day feature panel. The regularization
parameter was selected on validation. &
Compared top-\(k\) choices and score-based weights; the selected branch used a
score-weighted portfolio. &
Validation portfolio excess, not in-sample MSE. \\
\midrule
Stage 1 XGBoost &
XGBoost on the enhanced factor set. I searched tree depth/regularization
profiles and used recency weighting with a 20-day half-life in the best branch. &
Compared equal, rank, softmax, score, squared-score, and inverse-volatility
weighting with \(k\in\{30,35,40,50\}\). &
Validation excess return and positive excess rate. \\
\midrule
Base MLP &
MLP regressor with standardized features, ReLU activation, Adam optimizer, early
stopping, and \(L_2\) regularization. &
Top-\(k\) portfolio with capped score/softmax weighting. &
Validation excess return. \\
\midrule
Tuned MLP &
Same MLP family, with additional architecture and regularization search. This
included compact networks such as \((64,32)\) and wider networks such as
\((96,48)\). &
Portfolio search over top-\(k\) and softmax/score transformations. &
Canonical and robust self-test performance. \\
\midrule
LSTM &
Sequence model using recent 20-day feature histories for each stock. It was
trained to predict 3-day forward return from temporal input rather than static
rolling summaries. &
The best LSTM branch used top-35 with squared-score weighting. &
Canonical held-out excess return. \\
\midrule
Style-dynamic MLP &
MLP with market-state-aware portfolio behavior. It used market trend and stress
features to vary the allocation rule. &
Dynamic top-\(k\) and weighting rules depending on market state. &
Robust split and canonical split jointly. \\
\midrule
Recent-window MLP &
The final submitted model. It is a scikit-learn MLP regressor wrapped in a
standard scaler, trained only on a selected recent lookback window rather than
always using full history. Candidate lookbacks were \(63,126,189,252\), and
full history. &
Top-\(k\) selection with capped softmax weights. In the canonical self-test the
selected setup was lookback 189, \((96,48)\), top 30, and
\texttt{softmax\_2.0}. After the official data update, the final refreshed
submission selected lookback 189, the lighter \((64,32)\)-style MLP, top 35,
and \texttt{softmax\_2.0}. &
Primary: validation mean excess return. Secondary: positive excess rate and
rank IC. \\
\midrule
Router, ensemble, and excess-target variants &
Additional branches combined recent-window and style-dynamic experts, averaged
recent-window models, or changed the target to excess return with denoising and
recency weights. &
Either switched between experts, ensembled scores, or changed only the
portfolio rule while keeping one alpha model. &
Canonical, robust, and walk-forward consistency. \\
\bottomrule
\end{longtable}
}

The final model is intentionally not the most complex model. It is the model
whose training distribution matched the 3-day task best while still satisfying
the portfolio constraints and passing multiple held-out checks.

\subsection{Final Prediction Model}

The submitted model is a scikit-learn MLP regressor wrapped in a standard
scaler. The candidate architectures are small MLPs:
\[
    (64,32),\quad (96,48),
\]
with ReLU activation, Adam optimization, early stopping, and \(L_2\)
regularization. I intentionally kept the neural network small because the
training sample is a panel of only about one year of daily CSI500 data, and a
larger network was more likely to fit noise than to improve out-of-sample
portfolio returns.

The key modeling decision is to restrict the supervised training sample to a
recent rolling window. I search over
\[
    \{63,126,189,252,\text{full}\}
\]
trading-day lookbacks in the main recent-window experiments. The intuition is
that a 3-day prediction problem is sensitive to current market regime, so older
observations may dilute the signal. This hypothesis was supported by the
results: recent-window training improved the canonical held-out excess return
from \(+1.744\%\) for the tuned full-history MLP to \(+2.526\%\), and improved
the robust-split result from \(+0.684\%\) to \(+0.804\%\).

The final model is selected through validation, not manually fixed. In the
canonical self-test, the selected configuration was
\[
    \text{lookback}=189,\quad
    \text{architecture}=(96,48),\quad
    \text{top-}k=30,\quad
    \text{weight method}=\text{softmax\_2.0}.
\]
For the final refreshed submission generated after the official incremental
data update, the selected configuration was
\[
    \text{lookback}=189,\quad
    \text{architecture}=(64,32)\text{ with lighter regularization},\quad
    \text{top-}k=35,\quad
    \text{weight method}=\text{softmax\_2.0}.
\]
This demonstrates that \(k\) is not fixed at 30 by the rules or by my code; it
is selected from the candidate set. The rule is at least 30 positive-weight
stocks. The final refreshed CSV contains 35 stocks, sums to 1, and has maximum
single-name weight 10\%.

\subsection{Portfolio Weights}

Given predicted scores \(s_{i,t}\), I select the top \(k\) stocks and convert
scores to positive weights. The best final method uses a softmax transform:
\[
    \tilde{s}_{i,t} =
    \frac{s_{i,t}-\bar{s}_t}{\operatorname{sd}(s_t)+\epsilon},
    \qquad
    u_{i,t} = \exp(\tau \tilde{s}_{i,t}),
    \qquad
    w_{i,t} = \frac{u_{i,t}}{\sum_{j\in S_t} u_{j,t}},
\]
where \(S_t\) is the selected top-\(k\) set and \(\tau\) is a temperature
parameter. After this, weights are capped at 10\% and redistributed
proportionally until all constraints are satisfied. I also tested
volatility-adjusted softmax variants, but the main winning configuration used
\texttt{softmax\_2.0}, which placed more weight on the strongest predicted
names while respecting the 10\% cap.

\section{Experimental Design}

\subsection{Splits and Leakage Control}

All model comparisons use chronological splits. For a prediction date
\(t\), training labels use \(P_{t+3}\), so I cap the training pool at
\(t-3\) trading days. I also apply an embargo of 5 trading days between train
and validation or test blocks. This reduces overlap between forward-return
labels and avoids a hidden leakage channel in short-horizon returns.

The main canonical self-test uses:
\[
\begin{array}{ll}
\text{train} & \text{up to 2026-03-09},\\
\text{validation} & \text{2026-03-17 to 2026-03-30},\\
\text{test} & \text{2026-04-08 to 2026-04-21},\\
\text{embargo} & 5\text{ trading days}.
\end{array}
\]
I also use a stricter robust split with a 15-day validation block and a 20-day
test block:
\[
\begin{array}{ll}
\text{train} & \text{up to 2026-02-06},\\
\text{validation} & \text{2026-02-24 to 2026-03-16},\\
\text{test} & \text{2026-03-24 to 2026-04-21}.
\end{array}
\]
Finally, I run a 3-window walk-forward check to see whether a method survives
multiple historical test windows rather than only one favorable split.

\subsection{Evaluation Criteria and Hyperparameter Search}

For every branch, I selected hyperparameters using portfolio-level validation
metrics rather than regression loss alone. The main validation objective was
mean 3-day excess return after applying the actual portfolio construction rule.
I also tracked positive excess rate and rank IC to avoid selecting a model that
won only because of one lucky date.

The hyperparameter searches differed by model family. Ridge mainly searched
regularization and weighting. XGBoost searched tree shape, regularization,
recency weighting, top-\(k\), and score-to-weight transforms. MLP searched
architecture, regularization, top-\(k\), softmax temperature, and finally the
training lookback. The final model was chosen because it delivered the best
balance of canonical held-out performance and robustness, not because it was
the most complex model.

\section{Results}

\subsection{Main Model Comparison}

Table \ref{tab:main-results} reports the canonical held-out result for each
major attempt. The metric is mean 3-day excess return on the chronological test
block after portfolio construction. This is the relevant metric because the
competition evaluates the final constrained portfolio, not the model's
regression loss.

\begin{table}[h]
\centering
\caption{Canonical held-out Stage 1 results.}
\label{tab:main-results}
\begin{tabular}{p{0.27\textwidth}rrp{0.27\textwidth}}
\toprule
Model & Test excess & Gain vs. baseline & Interpretation \\
\midrule
3-day XGBoost baseline & \(+0.767\%\) & -- & first tradable benchmark \\
Ridge enhanced & \(+1.273\%\) & \(+0.506\%\) & feature engineering helped \\
Stage 1 XGBoost & \(+1.394\%\) & \(+0.627\%\) & nonlinear trees helped modestly \\
Stage 1 MLP & \(+1.741\%\) & \(+0.974\%\) & neural cross-section improved ranking \\
Tuned MLP & \(+1.744\%\) & \(+0.977\%\) & architecture tuning gave little extra gain \\
LSTM & \(+0.890\%\) & \(+0.123\%\) & sequence model underperformed \\
Style-dynamic MLP & \(+0.308\%\) & \(-0.459\%\) & robust in one split, weak canonical \\
\textbf{Recent-window MLP} & \(\mathbf{+2.526\%}\) & \(\mathbf{+1.759\%}\) & final Stage 1 model \\
\bottomrule
\end{tabular}
\end{table}

\textbf{Baseline XGBoost.}
The baseline produced \(+0.767\%\) canonical excess return. This result is
important because it establishes that the raw price-volume data contains some
short-horizon signal, but the portfolio still leaves a large amount of
performance on the table.

\textbf{Ridge enhanced.}
The enhanced Ridge model improved the canonical result to \(+1.273\%\). Since
Ridge is a linear model, this gain mainly came from better factor construction:
market-relative returns, range and gap features, downside volatility, beta,
idiosyncratic momentum, and cross-sectional ranks. This showed that factor
engineering mattered before adding model complexity.

\textbf{Stage 1 XGBoost.}
XGBoost on the enhanced factor panel reached \(+1.394\%\). This was better than
Ridge, but the improvement was moderate. My interpretation is that nonlinear
interactions among the factors were useful, but tree complexity alone was not
the main bottleneck.

\textbf{Base and tuned MLP.}
The first MLP improved the canonical result to \(+1.741\%\), and the tuned MLP
reached \(+1.744\%\). The jump from XGBoost to MLP suggests that standardized
continuous factors and smooth nonlinear interactions were useful for
cross-sectional ranking. However, the tiny gain from base MLP to tuned MLP
also showed that simply changing hidden-layer sizes was not enough for a large
step improvement.

\textbf{LSTM.}
The LSTM achieved \(+0.890\%\), above the baseline but below Ridge, XGBoost,
and MLP. This result suggests that explicit sequence modeling was not worth
the extra complexity for this dataset. The panel is short, the target is noisy,
and the rolling technical features already summarize much of the recent path.

\textbf{Style-dynamic MLP.}
The style-dynamic branch achieved only \(+0.308\%\) on the canonical split,
although it later showed a strong robust-split result. I did not submit it
because the canonical failure indicated that the dynamic rule was unstable
across market windows.

\textbf{Recent-window MLP.}
The recent-window MLP achieved \(+2.526\%\), the strongest canonical held-out
result. It improved over the tuned MLP by \(0.782\) percentage points and over
the Stage 1 XGBoost model by \(1.132\) percentage points. Because the model
family and portfolio rule were close to the tuned MLP, this is strong evidence
that the training-window choice was the main improvement.

\subsection{Robustness Checks}

\begin{table}[h]
\centering
\caption{Robustness tests for the strongest submission candidates.}
\begin{tabular}{lrrr}
\toprule
Model or evaluation & Test excess & Positive excess rate & Rank IC \\
\midrule
Tuned MLP robust split & \(+0.684\%\) & \(70.0\%\) & \(0.0610\) \\
Style-dynamic MLP robust split & \(+1.130\%\) & \(70.0\%\) & \(0.0653\) \\
\textbf{Recent-window MLP robust split} & \(\mathbf{+0.804\%}\) & \(60.0\%\) & \(0.0562\) \\
\textbf{Recent-window MLP walk-forward} & \(\mathbf{+0.783\%}\) & \(66.7\%\) & \(-0.0251\) \\
\bottomrule
\end{tabular}
\end{table}

The style-dynamic branch had the strongest robust-only result, but it performed
poorly on the canonical split. Therefore I did not use it as the final
submission model. The recent-window MLP was more balanced: it was the best
canonical model, beat the tuned MLP under the robust split, and remained
positive under multi-window walk-forward evaluation.

\subsection{Recent-Window Ablation}

\begin{table}[h]
\centering
\caption{Effect of recent-window training relative to full-history tuned MLP.}
\begin{tabular}{lrr}
\toprule
Evaluation & Tuned full-history MLP & Recent-window MLP \\
\midrule
Canonical held-out test excess & \(+1.744\%\) & \(\mathbf{+2.526\%}\) \\
Robust 20-day test excess & \(+0.684\%\) & \(\mathbf{+0.804\%}\) \\
\bottomrule
\end{tabular}
\end{table}

This is the cleanest evidence in the project. The model family and portfolio
rules are almost unchanged; the major change is restricting training to recent
market history. The improvement supports the hypothesis that the 3-day horizon
is regime-sensitive and that older data can be harmful.

\subsection{Failed or Non-Submitted Experiments}

\begin{table}[h]
\centering
\caption{Additional research branches that were not selected for submission.}
\begin{tabular}{p{0.27\textwidth}rrp{0.27\textwidth}}
\toprule
Branch & Canonical excess & Robust excess & Reason not submitted \\
\midrule
Excess-target weighted MLP & \(+0.993\%\) & \(-0.522\%\) & robust result became negative \\
Model-switch router & \(+1.842\%\) & \(+0.860\%\) & small gain, high complexity, weak router objective \\
Recent-window ensemble & \(+1.244\%\) & \(+0.638\%\) & below single recent-window model \\
Optimized portfolio layer & \(+1.424\%\) & -- & did not beat recent-window MLP \\
\bottomrule
\end{tabular}
\end{table}

These negative results are important. They show that simply adding more moving
parts did not reliably improve the submission metric. The best model was the
simpler recent-window MLP, which changed the training sample in a way that
matched the short-horizon nature of the task.

\textbf{Excess-target weighted MLP.}
This branch changed the supervised target from raw 3-day return to 3-day excess
return and added target denoising such as winsorization or z-scoring. It was
designed to align training more directly with the evaluation metric. The result
was weaker: \(+0.993\%\) canonical excess, \(-0.522\%\) robust excess, and
\(+0.233\%\) walk-forward excess. I interpret this as evidence that
cross-sectional raw return still contained useful stock-selection signal that
was partially removed by the excess-target transformation.

\textbf{Model-switch router.}
The router trained separate recent-window and style-dynamic experts and then
used market features such as recent index return, volatility, drawdown, breadth,
and trend to switch between them. It reached \(+1.842\%\) canonical excess and
\(+0.860\%\) robust excess, but it was expensive to train and its walk-forward
router objective was weak. I did not select it because the extra complexity was
not justified relative to the simpler recent-window MLP.

\textbf{Recent-window ensemble.}
The ensemble averaged or combined recent-window variants to reduce model
variance. It reached \(+1.244\%\) canonical excess and \(+0.638\%\) robust
excess. This was positive but below the single recent-window MLP, so ensembling
did not solve the main problem.

\textbf{Optimized portfolio layer.}
This branch focused on changing the score-to-weight rule while keeping the
alpha model similar. It achieved \(+1.424\%\) canonical excess. The result was
better than the baseline but below the recent-window model, so portfolio
optimization alone was not the highest-impact lever.

\section{Analysis}

\subsection{What Worked}

The largest improvement came from using a recent training window. For a 3-day
horizon, the predictive relation between technical factors and future returns
appears unstable across regimes. The recent-window model reduced exposure to
older regimes while preserving enough observations for MLP training. This
produced a large canonical gain and still survived the robust and walk-forward
checks.

The enhanced factor set also mattered. Market-relative return features helped
the model focus on stock-specific strength instead of simply selecting high
beta names in a rising market. Cross-sectional rank features matched the
portfolio construction problem, because the final task is to choose relative
winners among the same daily universe. Liquidity, volatility, beta, and
idiosyncratic momentum features gave the model information about whether a
short-term move was broad, stock-specific, or risk-driven.

The softmax portfolio construction was also effective. A uniform top-30
portfolio would ignore score strength, while an unconstrained allocation would
violate the 10\% cap or become too concentrated. Softmax weighting with a cap
provided a middle ground: it emphasized high-conviction names while remaining
valid under the competition rules.

\subsection{What Did Not Work}

More complex modeling was not automatically better. LSTM did not beat the tuned
MLP, suggesting that the available panel length was not enough for a sequence
model to dominate simpler cross-sectional models. The style-dynamic branch
improved the robust split but failed on the canonical split, so it was not safe
as a final submission model. The model-switch router was conceptually
reasonable but too complex relative to its gain. It slightly improved robust
excess return but reduced canonical performance and introduced a high
computational cost.

Changing the supervised target from raw 3-day return to excess 3-day return,
combined with denoising and light dynamic allocation, also did not help. That
branch produced \(+0.993\%\) canonical excess, \(-0.522\%\) robust excess, and
\(+0.233\%\) walk-forward excess. This suggests that the raw-return target was
capturing useful short-horizon stock selection signal that was weakened by
over-normalizing the label.

\subsection{Portfolio Interpretation}

The final refreshed Stage 1 submission after the official data update selected
35 names. The largest weights hit the 10\% cap, which is expected under
\texttt{softmax\_2.0}. The top holdings in the refreshed CSV were:

\begin{table}[h]
\centering
\caption{Top holdings in the refreshed Stage 1 submission.}
\begin{tabular}{lr}
\toprule
Stock code & Weight \\
\midrule
688037 & \(10.000\%\) \\
002281 & \(10.000\%\) \\
000062 & \(10.000\%\) \\
600208 & \(10.000\%\) \\
300136 & \(10.000\%\) \\
000539 & \(7.199\%\) \\
688166 & \(6.083\%\) \\
003031 & \(5.753\%\) \\
600487 & \(4.427\%\) \\
600884 & \(3.090\%\) \\
\bottomrule
\end{tabular}
\end{table}

I interpret these not as manually chosen ``favorite'' stocks, but as the
result of the model's short-horizon factor ranking. The cap prevents a few
names from dominating the portfolio completely. Because no news or fundamental
data was used, I do not claim that these holdings are justified by firm-level
events; the justification is strictly factor-based.

\section{Final Submission}

Before generating the final Stage 1 CSV, I updated the data according to the
competition instruction using the incremental download script. The final
submission was generated by
\[
    \texttt{mlp/day3\_stage1\_recent\_window/mlp\_model.py}.
\]
The refreshed submission metadata was:

\begin{itemize}
    \item prediction date: \texttt{2026-04-30};
    \item selected lookback: 189 trading days;
    \item selected training-window start: \texttt{2025-06-25};
    \item selected model: \texttt{narrow\_deep\_lighter\_reg};
    \item selected portfolio: top \(35\), \texttt{softmax\_2.0};
    \item validation rank IC: \(0.0531\);
    \item validation mean excess return: \(+1.881\%\);
    \item validation positive excess rate: \(80.0\%\).
\end{itemize}

The generated CSV satisfies all submission constraints: it has 35 positive
weights, total weight 1, no short positions, and maximum single-name weight
10\%.

\section{Conclusion}

The final Stage 1 system is a recent-window MLP with enhanced technical and
market-relative factors and capped softmax portfolio construction. The most
important empirical finding is that for a 3-day horizon, using a shorter recent
training window improved the submission-oriented metric more than adding model
complexity. The final model achieved \(+2.526\%\) canonical held-out excess
return, \(+0.804\%\) robust-split excess return, and \(+0.783\%\) 3-window
walk-forward excess return. Based on this balance of performance and
robustness, I selected \texttt{mlp/day3\_stage1\_recent\_window} as my final
Stage 1 model.

\appendix

\section{Implementation Artifacts}

\begin{itemize}
    \item Final model code:
    \texttt{mlp/day3\_stage1\_recent\_window/mlp\_model.py}.
    \item Feature code:
    \texttt{ridge/day3\_optimized/features.py}, imported through the MLP
    branch feature entrypoints.
    \item Canonical self-test:
    \texttt{submissions/mlp/day3\_stage1\_recent\_window/self\_test.json}.
    \item Robust self-test:
    \texttt{submissions/mlp/day3\_stage1\_recent\_window\_robustsplit/self\_test.json}.
    \item Walk-forward evaluation:
    \texttt{submissions/mlp/day3\_stage1\_recent\_window/walk\_forward.json}.
    \item Final refreshed submission:
    \texttt{submissions/stage1/submission\_stage1\_recent\_window\_updated\_20260503.csv}.
\end{itemize}

\section{Reproducibility Commands}

\begin{verbatim}
python download_data.py --update --end 20260503

.venv312/bin/python mlp/day3_stage1_recent_window/mlp_model.py \
  --out submissions/stage1/submission_stage1_recent_window_updated_20260503.csv \
  --json-out submissions/stage1/dev_stage1_recent_window_updated_20260503.json

.venv312/bin/python baseline/validate_submission.py \
  submissions/stage1/submission_stage1_recent_window_updated_20260503.csv
\end{verbatim}

\end{document}
